% !TeX root = ../../surface-lagrangian-torus-fibration.tex
% 2026-11-22 / v1.2: Source 保持下の成分別面積配分・finite cuts・actual 微小実現．
\section{Critical actions を直線に保つ相対 neck 接着}
\label{ng8:section}

葉層の相対 isotopy による補間だけでは，critical action curve の直線性は保たれない．
ここでは，共通の Hamiltonian circle と三重点の toric reference circles を使って
作用の加法定数を固定し，critical regions の面積を指定して補間する．
輸送を行うのは常に正の距離を置いた切断面であり，
$\mu=0,t=0$ の固定点集合を通る parallel transport は用いない．

\subsection{作用の基準と固定 moment 上の輸送}

二重曲線 $C=C_{ij}$ から三重点の小円板を除いた曲面を $P$ とする．
節~\ref{sec:necks} の actual neck には，周期 $2\pi$ の Hamiltonian circle と
正規化された moment $\mu$ がある．$i$ 側で normal moment が $\mu>0$，
$j$ 側で $-\mu>0$ となるように向きを取る．
有限個の nodal patches より十分薄い collar を選ぶ．
各 $t\ne0$ の neck では circle action は free であるが，
以下で primitive を取る全空間の tube には，中央 fiber の固定集合 $C^\circ$ も含める．

buffer を持つ全空間の tube は，$P$ と同じ homotopy 型の open curve へ retract するので，
$H^2=0$ である．
従って $\Omega=\dd\lambda_{\rm neck}$ と書ける．
circle averaging により $\lambda_{\rm neck}$ を invariant に取ると，
Cartan の公式と $\iota_V\Omega=-\dd\mu$ から
\[
 \dd\bigl(\lambda_{\rm neck}(V)\bigr)
 =-\iota_V\Omega=\dd\mu
\]
を得る．固定集合 $C^\circ$ 上では $V=0$ および $\mu=0$ なので，
$\lambda_{\rm neck}(V)=\mu$ となる．
以下の作用は，この一つの primitive で正規化する．
これは各 puncture の toric action の加法定数を独立に零にする操作ではない．

各 puncture end の toric chart で，circle generator $\alpha$ ともう一つの
primitive generator $\beta_\ell$ を固定する．
normal moment と残る toric moment の affine cuts で
compact な reduced surface $P_{t,\mu}$ を切り出す．
その境界 circle の第二作用を $R_\ell(\mu)$ と書く．
各 end の toric moment を $\nu_\ell$ とすると，これは
$\nu_\ell+c_\ell^{\rm ref}$ の cut への制限である．
$c_\ell^{\rm ref}$ は $\lambda_{\rm neck}$ で決まる定数で，
両側の成分と全ての $t$ に共通である．
実際，$\beta_\ell$ の toric flow を $R_\theta$，generator を $V_{\beta_\ell}$ とすると，
その orbit action は
\[
 \widetilde R_\ell=
 \frac1{2\pi}\int_0^{2\pi}R_\theta^*
       \bigl(\lambda_{\rm neck}(V_{\beta_\ell})\bigr)\,\dd\theta
\]
である．Cartan の公式の $\mathcal L_{V_{\beta_\ell}}\lambda_{\rm neck}$ 項は
一周期で積分すると零になるので，
$\dd\widetilde R_\ell=-\iota_{V_{\beta_\ell}}\Omega=\dd\nu_\ell$ を得る．
従って同じ連結な全空間の toric end で
$\widetilde R_\ell=\nu_\ell+c_\ell^{\rm ref}$ である．
この計算に $\lambda_{\rm neck}$ 自体の全 $T^2$-invariance は必要ない．
従って $R_\ell(\mu)$ は実数係数の affine function であり，$t$ に依存しない．
cut 自体の勾配が整数であることは要求しない．

boundary circle から同じ regular Reeb arm 上の circle $\gamma$ までの
reduced annulus を $A$ とすれば，向きと lift を一貫して選ぶことで
\begin{equation}\label{ng8:relative-action}
 J_{\beta_\ell}(\mu,\gamma)-R_\ell(\mu)
 =\frac1{2\pi}\int_A\omega_{t,\mu}
\end{equation}
となる．ここで $J_\alpha=\mu$ と正規化する．
lift の整数変更は $J_\beta$ に整数倍の $\mu$ を加える．
critical circle への値は regular annuli の exhaustion により定める．
内部 Reeb edge の作用差も，同じ式をその annulus に適用して定める．

\begin{lemma}[切断面における作用の加法定数の保存]
\label{ng8:transport-intercept}
固定した $0<\delta$ と $\mu=\delta$ または $\mu=-\delta$ の近傍で，
成分側から nearby fiber への輸送は，
toric reference circles で正規化した作用の加法定数を保存する．
従って，成分側の critical action germ が
$J_\beta=c+p\mu$ なら，輸送後も同じ $c,p$ を持つ．
\end{lemma}

\begin{proof}
補題~\ref{glue:transport} の horizontal vector field は circle invariant で，
全空間上の $\mu$ を保存する．固定 $\mu\ne0$ で symplectic reduction を取ると，
その vector field は reduced symplectic forms を保つ horizontal field に降りる．
三重点の端では二つの toric moments を保存するので，
reference circle とその作用 $R_\ell(\mu)$ も保存される．
式~\eqref{ng8:relative-action} の annulus を運べば，その面積が保存される．
critical limits に対しても exhaustion により同じ結論が従う．

一つの全空間 primitive を使うと，作用そのものの保存も直接確認できる．
輸送 path を $\tau\mapsto\tau t$，その horizontal lift を $W_\tau$，
transport を $\mathcal P_{\tau t}$ とする．
固定した $\mu=\pm\delta$ の buffers を十分小さく取れば，全輸送軌跡は上の tube に入る．
$W_\tau$ は fiber の tangent space と $\Omega$-orthogonal なので，
Cartan の公式を fiber tangent に制限すると
\[
 \left.\mathcal L_{W_\tau}\lambda_{\rm neck}\right|_{TX_{\tau t}}
 =\dd_{X_{\tau t}}\bigl(\lambda_{\rm neck}(W_\tau)\bigr).
\]
従って，輸送の定義域上で
\begin{equation}\label{ng8:exact-transport}
 \mathcal P_t^*(\lambda_{\rm neck}|_{X_t})
   -\lambda_{\rm neck}|_{D_i}
 =\dd\!\left(\int_0^1
       \mathcal P_{\tau t}^*\bigl(\lambda_{\rm neck}(W_\tau)\bigr)
       \,\dd\tau\right).
\end{equation}
閉じた transported cycles の積分は正確に保存される．
この等式は primitive の period を任意に再正規化する操作ではない．

切断面の小さい open moment interval 上でこの議論を行うので，
作用微分だけでなく関数そのものが保存される．
定理の証明に必要なのは固定した二つの切断面とその buffers だけである．
$\delta\to0$ に一様な輸送を要求していない．
\end{proof}

\begin{lemma}[相対 reduced area の affine 性]
\label{ng8:relative-dh}
固定した十分小さい $t\ne0$ に対し，
上記の toric moment cuts で切り出した $P_{t,\mu}$ の全 reduced area は
$\mu$ の affine function である．
その値は，両端の切断面で中央成分から輸送した reduced area と一致する．
\end{lemma}

\begin{proof}
$t\ne0$ の neck に固定点はなく，moment interval 上の reduced surfaces は
境界を含めて同じ曲面 $P$ と同定できる．
principal circle bundle に connection を選ぶと，reduced cohomology class の変化は
その curvature class に比例する．これは Duistermaat--Heckman の局所的な計算である
\cite[Chapter~10]{DHJH}．境界付き曲面では絶対 cohomology class だけから総面積は
決まらないため，ここで必要な主張は次の Stokes の計算で直接証明する．

境界項も確認する．$P$ は境界を持つので circle bundle は自明化できる．
この自明化と各 toric end の framing の差は，固定した整数 winding である．
固定 $\mu$ 上で $\lambda_{\rm neck}-\mu\dd\theta$ は reduced surface の primitive に降りる．
従ってその境界 circle の積分は $R_\ell(\mu)+m_\ell\mu$，
$m_\ell\in\ZZ$，という affine function になる．
Stokes の定理により全 reduced area は，これらの境界作用の符号付き和の
$2\pi$ 倍であり，affine である．
moving cut を用いる場合の境界の寄与も $R_\ell(\mu)$ に既に含まれている．
両端の値の一致は補題~\ref{ng8:transport-intercept} の reduced transport による．
\end{proof}

\begin{lemma}[固定した fiber の free band における相対面積]
\label{ng8:fixed-fiber-relative-dh}
一つの $t\ne0$ と compact moment interval $J$ を固定する．
Actual $(E_t,\omega_t)$ 上に周期 $2\pi$ の自由な Hamiltonian circle action があり，
$\iota_V\omega_t=-\dd\mu$，$E_t/S^1\simeq J\times P$ とする．
ここで $P$ は非空の境界を持つ compact connected oriented surface である．
作用不変な primitive $\lambda_t$ が
$\dd\lambda_t=\omega_t$，$\lambda_t(V)=\mu$ を満たすと仮定する．
各 boundary component の toric framing と cut を固定し，
境界の向きに合わせた第二作用 $R_{\ell,t}(\mu)$ が affine であるとする．
このとき，$\mu$ に依存しない整数 $k_{\ell,t}$ があって
\begin{equation}\label{ng8:eq:fixed-fiber-area}
 \int_{P_{t,\mu}}\omega_{t,\mu}
 =2\pi\sum_\ell\bigl(R_{\ell,t}(\mu)+k_{\ell,t}\mu\bigr)
\end{equation}
となる．特に total reduced area は $J$ 上で affine である．
Reference constants の $t$ 依存を許し，$t=0$ を含む全空間への形式の延長，
central component からの輸送，又は $t$ に一様な評価は仮定しない．
\end{lemma}
\begin{proof}
$H^2(J\times P;\ZZ)=0$ なので principal circle bundle を自明化し，
周期 $2\pi$ の angular coordinate $\theta$ を取る．
$\lambda_t-\mu\dd\theta$ は horizontal かつ invariant であり，
固定 $\mu$ の reduction 上の primitive $\beta_{t,\mu}$ に降りる．
この自明化の boundary section と toric framing の差は整数 winding
$k_{\ell,t}$ で，連結な $J$ 上で一定である．従って
$\int_{\partial_\ell P}\beta_{t,\mu}
=2\pi(R_{\ell,t}(\mu)+k_{\ell,t}\mu)$ であり，Stokes の定理で結論を得る．
Toric moment $\nu_\ell$ を用いる場合にも，上の orbit-average 計算により
第二作用は $\nu_\ell+c_{\ell,t}^{\rm ref}$ である．
この定数は固定 $t$ の連結な end で一定であれば足り，全 $t$ で同じ値である必要はない．
\end{proof}

\begin{example}[Source 固定点を横切るときの面積の折れ]
\label{ng8:source-dh-wall}
$\CC^2$ 上で $I_k=|z_k|^2/2$，角度の周期を $2\pi$ とし，
標準形式 $\omega=\sum_{k=1}^2\dd I_k\wedge\dd\theta_k$，
重み $(1,-1)$ の circle action，$\mu=I_1-I_2$ を取る．
固定した $c>0$ に対し $I_1+I_2\le c$ を reduction すると，$|\mu|<c$ でその面積は
\begin{equation}\label{ng8:eq:source-reduced-area}
 \operatorname{Area}(P_\mu)=\pi(c-|\mu|)
\end{equation}
となる．$\mu=0$ では singular reduction の regular part 上の積分を面積とし，
これは両側からの連続極限にも一致する．
\end{example}
\begin{proof}
両座標が非零の部分で $I=I_2$，$\phi=\theta_1+\theta_2$ と置くと，
$\phi$ の周期は $2\pi$ であり，reduced form は $\dd I\wedge\dd\phi$ である．
許される interval は
\[
 \max\{0,-\mu\}\le I\le(c-\mu)/2.
\]
下端は $\mu\ne0$ では滑らかな reduction の一点，$\mu=0$ では固定点の像であり，
いずれも上記の積分に面積を加えない．従って面積は
$2\pi\{(c-\mu)/2-\max(0,-\mu)\}=\pi(c-|\mu|)$ である．
外側の toric boundary action $I_2=(c-\mu)/2$ は affine だが，
原点は $\mu=0$ の固定点であり，全 interval に対する free-band の自明化は存在しない．
実際，$\mu>0$ では $\theta_1=0$，$\mu<0$ では $\theta_2=0$ の section を使うと，
boundary winding はそれぞれ $0$ と $1$ となり，上の二つの affine pieces を与える．
\end{proof}

\begin{remark}[Actual KE の相対 atlas へ使う際の入力]
\label{ng8:fixed-fiber-area-scope}
前例は，free-band の面積式を source 固定点を横切る interval へ拡張できないことを示す．
Actual P0 全体の reduced area を計算した例ではなく，既存の atlas 存在定理や
P0 の訂正でもない．Actual KE 形式に $t=0$ までの smooth total-space extension を
仮定することなく，前補題を各固定 $t$ の free compact band に適用できる．
P0 の周期 $1$ の規約からは，生成場と canonical moment をともに $2\pi$ で割って用いる．

ただし補題\ref{ng8:parametric-moser}を続けて使うには，marked identifications と
端の open foliation germs，正の target region areas，その総量と actual reduced area の一致，
および\eqref{ng8:protected-area-margin}を別に確認する．
固定領域 $K$ には保存する全 collars と caps を含め，同補題の相対構成の条件を満たす必要がある．
同じ normal moment と end の葉を保つことだけでは，これらの面積条件は従わない．
Source の近傍を保持する場合は変更の支持をそこから離し，
critical-area の調整と core collar での一致を別々に確保する．
保持 annuli により $P\setminus K$ が分かれる場合は，
定理\ref{ng8:componentwise-critical-area}の成分別の配分条件を用いる．
系\ref{ng8:p0-small-area-flexibility}は source の全飽和近傍を保持した
微小変更を与えるが，任意の所定 target への到達を主張しない．
固定 compact と cuts を選んだ後に $t$ の閾値を取る構成でも，
この相対実現の条件そのものを対角選択で代用することはできない．
\end{remark}

\subsection{Critical regions の面積を保つ補間}

球面上の固定された marked Morse foliation の critical leaves は，
有限個の puncture basins と，内部 Reeb edges に対応する annuli に球面を分ける．
以下では，これらをまとめて critical regions と呼ぶ．
各 critical leaf 自体の面積は零である．
境界付き曲面を用いる場合は，puncture caps を補助的に付けて同じ記述を用いてよい．
caps と保存する puncture collars の和を，互いに素な固定 disks の内部に取る．
この固定領域を $K$ と書く．その上の foliation と面積形式の parameter family は
あらかじめ与えておき，以下の操作では変更しない．
面積形式が parameter に依存しないことを要求しているわけではない．

\begin{lemma}[Parameter に相対的な critical-area 指定]
\label{ng8:parametric-moser}
$u\in[u_-,u_+]$ に対して，球面上の正の面積形式 $\omega_u$ と，
固定 puncture collars を持つ marked Reeb tree を与える．
両端の open intervals には，同じ相対 isotopy class に属する
Morse foliation germs が指定されているとする．

各 critical region $D_h$ の target area $a_h(u)>0$ が滑らかで，
その和が $\int_C\omega_u$ に等しく，変更しない全領域の面積に対して
\begin{equation}\label{ng8:protected-area-margin}
 a_h(u)>\int_{D_h\cap K}\omega_u
\end{equation}
を満たすとする．$K$ には caps だけでなく，保存する collars 全体を含める．
両端では指定された germs がこの面積を持つと仮定する．
このとき，端の germs と collars を保ち，
全 $u$ で $\operatorname{area}_{\omega_u}(D_h)=a_h(u)$ を満たす
滑らかな marked foliation family を構成できる．
追加の compact parameters を用いる場合は，初期の marked interpolation と
region identifications がそれらに関しても滑らかに与えられていると仮定する．
その仮定の下で，面積指定と Moser の操作も parameter に相対的に行える．
\end{lemma}

\begin{proof}
まず補題~\ref{neck:marked-interpolation} により，marked foliations を補間する．
基準 foliation と滑らかな diffeomorphisms を使って各 region を同定すれば，
それぞれの面積は parameter に関して滑らかである．

各 region 内の，固定領域 $K$ を避けた部分に非負 bump functions を取る．
正の density を掛けて面積を指定する方法は補題~\ref{per:lobe-areas} と同じである．
\eqref{ng8:protected-area-margin} と parameter interval の compactness により，
少し大きい固定近傍でも一様な面積の余裕がある．
その近傍では density を一に保ち，それ以外でまず十分小さくした後，
各 region の内部に支持された bump を加える．
局所的には，region ごとの bump coefficient が面積の独立な座標となるので，
正値性を保つ滑らかな解を取れる．
同じ target integrals を持つ正の面積形式の集合は凸であるから，
parameter 空間上の partition of unity で局所解を貼れる．
両端では元の形式が既に条件を満たしているので，そこでは変更しない．
こうして，固定 collars 上で $\omega_u$ に一致し，
指定された region integrals を持つ正の形式 $\widehat\omega_u$ を得る．

両形式の差は $K$ の近傍で零であり，その総積分も零である．
固定 disks の補集合は連結なので，その interior の compactly supported
二次 de Rham cohomology は積分により $\RR$ と同定される．
従って差には，$K$ の近傍で零となる primitive がある．
同じことは，境界付き曲面の記述では $H^2(P,\partial P;\RR)\simeq\RR$ から従う．
固定した相対 de Rham 解作用素を用いて primitive を $u$ に滑らかに選び，
差が零である両端の intervals では primitive も零に取る．
凸結合に Moser の方法を適用して
\[
 \psi_u^*\widehat\omega_u=\omega_u
\]
を満たす diffeomorphisms を得る．
端の intervals および fixed collars では $\psi_u=\id$ にできる．
最初に補間した $u$ における foliation を $\psi_u^{-1}$ で移せば，
元の $\omega_u$ に関して所期の面積を持つ．
全ての操作は parameter に滑らかで，Morse 型と marked Reeb tree を保つ．
\end{proof}

\begin{proposition}[Straight critical actions を保つ actual neck]
\label{ng8:straight-neck}
actual Hamiltonian neck の両端の fibration germs を与える．
両 germ の reduced foliations は，同じ実際の reduced surface 上で
boundary collars を固定する $\operatorname{Diff}_0$ の同じ orbit に属するものとする．
従って同じ embedded marked Reeb tree を持ち，
共通の cycle framings と toric reference actions で書いた critical actions が
同じ affine functions
\[
 J_{\beta_\ell}^{\rm crit}=c_\ell+p_\ell\mu,\qquad p_\ell\in\ZZ
\]
の制限であると仮定する．
これらから式~\eqref{ng8:relative-action} で指定される
basin および internal-annulus areas が正であり，
\eqref{ng8:protected-area-margin}の固定領域に対する余裕を持つとする．
このとき，全 interval にわたり critical actions が上の affine functions となる
Lagrangian fibration を，両端の open germs を保って構成できる．
\end{proposition}

\begin{proof}
boundary reference actions は affine なので，
target region areas は $\mu$ の affine functions である．
その和と total reduced area は，補題~\ref{ng8:relative-dh} により共に affine で，
両端で一致する．従って interval 全体で一致する．
正の target areas には compact interval 上で一様な正の下限がある．
puncture cuts と保存する collars は最初に十分小さく選び，
各 basin の target area から固定領域の面積を引いた値に一様な正の余裕を持たせる．
内部 annulus の target area の正値性は，命題の仮定として保持する．

補題~\ref{ng8:parametric-moser} を reduced surfaces に適用する．
端で一致する marked foliation family を得て，
その circle lifts の connected-fiber quotient を取る．
補題~\ref{neck:lift} により正則 fibers は Lagrangian であり，
critical fibers は $S^1\times(S^1\vee S^1)$ となる．
式~\eqref{ng8:relative-action} により critical actions は所期の affine functions である．
三重点の端では toric collars を固定しているので，既存の triple blocks と一致する．
\end{proof}

\subsection{中央成分の面積計算から nearby fiber への判定}

\begin{corollary}[中央 fiber による十分条件]
\label{ng8:central-criterion}
両成分の reduced foliations が同じ実際の曲線上の共通の
$\operatorname{Diff}_0(P\operatorname{ rel\ collars})$-orbit に属するとする．
両成分の collar germs で critical actions が整数勾配の affine functions であり，
共通の neck framings へ移した後の勾配が一致するとする．
さらに $\mu=0$ における strict double curve 上で，
対応する critical regions の実際の面積が両側で一致するとする．
このとき，十分小さい nearby fiber で命題~\ref{ng8:straight-neck} の接着ができる．
\end{corollary}

\begin{proof}
中央成分上では，$\mu=0$ で normal circle が潰れ，
式~\eqref{ng8:relative-action} は strict double curve 上の area difference になる．
framing の整数変更による $m\mu$ はこの点で零である．
三重点の toric reference actions は共通に固定されているので，
実際の basin areas の一致は critical action functions の定数項の一致を与える．
内部 annuli の作用差も，同じ面積公式で一致する．
Reeb graph は tree なので，external reference から各 critical block をたどり，
block 内の cycle relations と内部 annulus の作用差を用いれば，
全 critical actions の定数項が共通に定まる．
この過程で現れる整数 framing changes は $\mu=0$ では寄与しない．
勾配の仮定と合わせ，両側は同じ affine functions となる．

中央 curve の各 critical region の面積は正である．
puncture cuts と固定 collars を先に十分小さく取り，その後で $\delta>0$ を小さくすれば，
対応する affine target areas は $[-\delta,\delta]$ 全体で正となり，
\eqref{ng8:protected-area-margin}も満たす．

次に固定した正の切断値 $\mu=\delta$ と負の切断値 $\mu=-\delta$ および
その open buffers だけを nearby fiber へ運ぶ．
補題~\ref{ng8:transport-intercept} が同じ affine functions を保存する．
これで有限 $t$ における加法定数も一致した．
残りの moment interval には命題~\ref{ng8:straight-neck} を適用する．

この議論は，中央 fiber の面積一致から有限 $t$ の一致を無条件に推論するものではない．
固定 moment 上の相対 symplectic transport と toric reference normalization が，
その二つの計算を接続している．
\end{proof}


\subsection{固定領域の補集合ごとの面積条件}
\label{ng8:componentwise-area}

元の補題\ref{ng8:parametric-moser}では，caps と puncture collars を含む
固定 disks の補集合が連結であることを使っている．ここでは，source の全葉の
近傍などを追加して固定し，補集合が非連結となる場合を扱う．
新しい fixed set の下で初期の marked interpolation を作れることは，
面積条件とは別の入力である．特に，抽象的な labelled Reeb tree の一致を
$\operatorname{Diff}_0(P\operatorname{ rel}K)$ による同定の代わりには用いない．

$P$ を compact oriented smooth surface とし，$K\subset P$ を smooth boundary を
持つ固定 compact codimension-zero submanifold とする．$\partial P\ne\varnothing$
の場合はその collar を $K$ に含める．
\[
 P\setminus K=\coprod_{\alpha=1}^q U_\alpha
\]
は有限個の connected open surfaces の和とする．$K$ の各成分が disk であるとは
仮定しない．$\operatorname{Diff}_0(P\operatorname{ rel}K)$ は，$K$ のある近傍で
恒等な isotopy によって恒等写像と結ばれる diffeomorphisms を表す．
従ってその各元は各 $U_\alpha$ をそれ自身へ写す．

Parameter space $B$ は compact smooth manifold with corners とする．
$B$ 上の finite positive measures $\nu_b$ と，smooth な初期 marked
foliation family $\mathcal F_b^0$ を与える．各 $U_\alpha$ 上では
$\nu_b=\omega_b$ は正の smooth area form であり，固定した自由な boundary collars
を含む $\overline U_\alpha$ 上で $b$ にも smooth と仮定する．
$K$ の内部の measure が smooth area form であること，及びその parameter smoothness
は仮定しない．$\nu_b(\partial K)=0$ とする．
Critical leaves と境界の和を $\Sigma_b$ と書き，その補集合の labelled critical
regions を $D_1(b),\ldots,D_H(b)$ とする．$\nu_b(\Sigma_b)=0$ と仮定する．
各 $b_0\in B$ の近傍で，$K$ と各 $U_\alpha$ を保ち，
$\Sigma_{b_0},D_h(b_0)$ を $\Sigma_b,D_h(b)$ へ写す smooth diffeomorphism family
が与えられていると仮定する．これは region と固定領域の incidence の
実際の smooth trivialization であって，その抽象的な個数の一致だけではない．
$K$ 上の measure 又は foliation が $b$ に依存しないことは要求しない．
変更後は，各 $b$ で与えられた $K$ 上のデータをそのまま保つ．

\[
 E=\{(h,\alpha):D_h(b)\cap U_\alpha\ne\varnothing\},\qquad
 \kappa_h(b)=\nu_b(D_h(b)\cap K),\qquad
 M_\alpha(b)=\int_{U_\alpha}\omega_b
\]
と置く．$D_h(b)$ と $U_\alpha$ は開なので，$E$ の条件は交わりの内部が
非空であることと同じである．上の trivializations により $E$ は固定され，
$M_\alpha$ と各 free cell の積分は smooth である．
$\kappa_h$ 自体の smoothness は以下で使わない．

\begin{lemma}[固定支持を持つ primitive の成分別構成]
\label{ng8:componentwise-primitive}
$\eta_b$ を $\coprod_\alpha U_\alpha$ 内の一つの compact set に支持された
smooth family of two-forms とする．$K$ の近傍で零となる smooth family of
one-forms $\beta_b$ で $d\beta_b=\eta_b$ を得るための必要十分条件は
\[
 \int_{U_\alpha}\eta_b=0\qquad(1\le\alpha\le q)
\]
である．Primitive は $\eta_b$ に線型に選べ，従って $\eta_b=0$ となる
parameter 領域では $\beta_b=0$ とできる．支持も一つの compact set に収まる．
\end{lemma}

\begin{proof}
必要性は compact support に対する Stokes の定理による．十分性は各成分内で
示せばよい．支持を覆う有限個の oriented coordinate rectangles と，それらを
一つの基準 rectangle へ結ぶ有限本の経路を選ぶ．経路も有限個の重なる rectangles
で覆い，partition of unity によって $\eta$ を局所 forms に分ける．
各局所 form から，その積分と同じ質量を持つ固定 bump two-form を引くと，
積分零の局所 form を得る．残った bump masses は，重なる rectangles 内の
積分一の bumps を順に用いて基準 rectangle へ運べる．総積分が零なので，
基準 rectangle に最後の質量は残らない．

ここで必要な局所 primitive を明記する．Rectangle 内で
$\eta=f(x,y)\,dx\wedge dy$，$\int f=0$ とし，$x$ 方向の compactly supported
function $\rho$ を $\int\rho=1$ となるよう固定する．
\[
 g(y)=\int_{\mathbb R} f(x,y)\,dx,\qquad
 F(x,y)=\int_{-\infty}^{x}\bigl(f(s,y)-\rho(s)g(y)\bigr)\,ds,
 \qquad
 G(y)=\int_{-\infty}^{y}g(s)\,ds.
\]
すると $F,G$ は compactly supported であり，
\[
 \beta=F\,dy-\rho G\,dx,
 \qquad d\beta=f\,dx\wedge dy.
\]
重なる bumps の差にも同じ式を用いる．全ての選択は固定でき，演算は線型な
積分・乗算と有限和だけである．従って smooth parameter dependence と固定支持を
同時に得る．各 $U_\alpha$ の結果を足し，$K$ の近傍へ零で延長する．
\end{proof}

\begin{theorem}[有限面積配分と parameter-relative な実現]
\label{ng8:componentwise-critical-area}
上の設定で target areas $a_h(b)>0$ を与え，
$r_h(b)=a_h(b)-\kappa_h(b)$ が smooth であると仮定する．
$a_h$ と $\kappa_h$ が個別に smooth である必要はない．
初期 family の各 slice と同じ $\operatorname{Diff}_0(P\operatorname{ rel}K)$-orbit
に属し，$K$ 上の全 foliation germs を保ち，
\[
 \nu_b(D_h^{\rm new}(b))=a_h(b)
\]
となる smooth marked foliation family が存在するための必要十分条件は，
各 $b$ に対して有限行列 $x_{h\alpha}(b)$ が存在し，
\begin{align}
 x_{h\alpha}&>0&&((h,\alpha)\in E),&
 x_{h\alpha}&=0&&((h,\alpha)\notin E),\label{ng8:transport-support}\\
 \sum_\alpha x_{h\alpha}&=a_h-\kappa_h&& (1\le h\le H),&
 \sum_h x_{h\alpha}&=M_\alpha&& (1\le\alpha\le q)
 \label{ng8:transport-margins}
\end{align}
を満たすことである．ここで存在は，初期 marked family とその smooth
trivializations が既に与えられた範囲での主張である．

より具体的には，この条件の下で smooth $\psi_b\in
\operatorname{Diff}_0(P\operatorname{ rel}K)$ と smooth positive densities
$\rho_b$ を構成でき，$\widehat\nu_b=\rho_b\nu_b$ と置くと
\[
 \psi_b^*\widehat\nu_b=\nu_b,
 \qquad \mathcal F_b^{\rm new}=\psi_b^{-1}(\mathcal F_b^0)
\]
となる．$\rho_b=1$ を $K\cup\Sigma_b$ の近傍に要求してもよい．
この最後の条件は density の条件であり，$\psi_b$ が $\Sigma_b$ を固定するという意味ではない．
Free part 上では $\widehat\omega_b=\rho_b\omega_b$ が smooth positive area forms である．

初期 family が指定された open end germs と一致し，そこで
$a_h(b)=\nu_b(D_h(b))$ であるなら，その germs を保って構成できる．
より一般に，この等式が成立する開集合 $O\subset B$ と
$O_0\Subset O$ を指定すれば，$O_0$ 上で $\psi_b=\mathrm{id}$ とできる．
追加の compact parameters についても上の smooth trivializations が与えられていれば，
同じ結論が同時に成立する．
\end{theorem}

\begin{proof}
\emph{必要性．}
ある $\psi\in\operatorname{Diff}_0(P\operatorname{ rel}K)$ が初期 marked partition
を新しいものへ移すとする．$\psi$ は各 $U_\alpha$ を保ち，$K$ 上では恒等なので，
各 $D_h^{\rm new}\cap U_\alpha$ の非空性は初期 cell の非空性と同じであり，
$D_h^{\rm new}\cap K=D_h\cap K$ である．
\[
 x_{h\alpha}=\int_{D_h^{\rm new}\cap U_\alpha}\omega
\]
と置けば，非空 open cell の面積は正であり，row と column の和は
\eqref{ng8:transport-margins}になる．
従って，この必要性は指定された relative orbit の中での主張である．
別の marked isotopy class の領域が同じ incidence を持つとは主張していない．

\emph{Pointwise feasibility から smooth な配分へ．}
非空 entry だけを並べ，row と column の和を取る固定線型写像を $A$ と書く．
右辺を $y(b)$ とする．ある $b_0$ で正の解 $x^0$ を取り，
$\operatorname{im}A$ 上の線型 right inverse $R$ を固定する．
Pointwise feasibility により $y(b)-y(b_0)\in\operatorname{im}A$ なので，
\[
 x(b)=x^0+R\bigl(y(b)-y(b_0)\bigr)
\]
は $b_0$ の近傍で smooth な正の解である．解集合は凸だから，有限 cover と
parameter partition of unity により global smooth positive solution を得る．
$O$ 上では actual cell areas $\int_{D_h(b)\cap U_\alpha}\omega_b$ が正の解なので，
これを $O_0$ の近傍で選べる．

\emph{正の density の構成．}
まず一つの小さい parameter neighborhood を取り，指定された trivialization で
全 regions と $K$ を固定する．各非空 cell に一つずつ compactly supported
nonnegative bump $\varphi_{h\alpha,b}$ を取り，
\[
 \int_{D_h(b)\cap U_\alpha}\varphi_{h\alpha,b}\,\omega_b=1
\]
と正規化する．各 bump は $K$ と critical leaves から離して取る．

$K$ の外側の十分薄い近傍と $\Sigma_b$ の十分薄い近傍の，どの cell における
面積も任意に小さくできる．前者は $U_\alpha$ 内で $K$ へ収束する部分の面積が
零へ向かうこと，後者は $\Sigma_b$ の面積が零であることによる．
Smooth trivialization と parameter compactness により，これらの選択は
局所 parameter neighborhood 上で一様に行える．
従って，$K\cup\Sigma_b$ の近傍で $1$，少し大きい近傍の外で小さい正数
$\varepsilon$ となる smooth positive function $r_b$ を取り，
\[
 m_{h\alpha}(b):=
 \int_{D_h(b)\cap U_\alpha}r_b\,\omega_b
 <x_{h\alpha}(b)
 \quad ((h,\alpha)\in E)
\]
とできる．この際 bumps の支持を $r_b=\varepsilon$ の領域に置ける．
\[
 \rho_b=r_b+
 \sum_{(h,\alpha)\in E}
 \bigl(x_{h\alpha}(b)-m_{h\alpha}(b)\bigr)\varphi_{h\alpha,b}
\]
は正で，$K\cup\Sigma_b$ の近傍では $1$，各 cell の積分は指定された
$x_{h\alpha}$ である．$\widehat\omega_b=\rho_b\omega_b$ と置く．
$K$ 上では density を $1$ に保つので，row 和と
$\kappa_h+r_h=a_h$ から $\widehat\nu_b(D_h(b))=a_h(b)$ となる．
Column 和から，各 $U_\alpha$ の総積分は $M_\alpha$ になる．

この構成を有限 parameter cover 上で行い，正の densities の凸性により貼る．
Constraints は線型なので全積分を保つ．$O$ 上には $\rho_b=1$ を局所解として
用い，partition of unity を $O_0$ の近傍でこの解だけになるよう選ぶ．
また有限個の近傍を縮めれば，$K$ の一つの共通近傍上で全ての densities が
$1$ になる．従って $\widehat\omega_b-\omega_b$ の支持は
$P\setminus K$ の一つの compact set に収まる．

\emph{Componentwise relative Moser．}
Column 和により
\[
 \int_{U_\alpha}(\widehat\omega_b-\omega_b)=0
\]
なので，前補題から smooth $\beta_b$ を
$d\beta_b=\widehat\omega_b-\omega_b$，$\beta_b=0$ near $K$ と選べる．
$O_0$ 上では $\beta_b=0$ とする．
Free part 上の $\omega_{b,s}=(1-s)\omega_b+s\widehat\omega_b$ は正なので，
\[
 \iota_{X_{b,s}}\omega_{b,s}=-\beta_b
\]
は smooth vector fields を定める．これらは $K$ の共通近傍で零であり，
compactness により $0\le s\le1$ の全 flow が存在する．
その time-one map $\psi_b$ は smooth で，
free part 上で $\psi_b^*\widehat\omega_b=\omega_b$ となる．
$K$ 上では恒等なので，measure として全 $P$ 上で
$\psi_b^*\widehat\nu_b=\nu_b$ となる．
Isotopy 全体が $K$ の近傍で恒等なので relative marking を保つ．
また $O_0$ 上では全 flow が恒等である．
従って $D_h^{\rm new}(b)=\psi_b^{-1}(D_h(b))$ は所望の面積を持つ．
Morse 型，full critical leaf の型，region labels は diffeomorphism で保存される．
\end{proof}

\begin{example}[全体の総量と protected-area margin だけでは足りない]
\label{ng8:annular-protection-obstruction}
球面上に四つの puncture basins $D_1,\ldots,D_4$ と一つの internal annulus を持つ
marked Morse foliation を固定し，正の面積形式 $\omega$ を固定する．
各 puncture の小さい cap disk を保護する．さらに $D_1$ の regular circle の
小さい飽和 annulus $A\Subset D_1$ を固定し，これらの和を $K$ とする．
$A$ と $D_1$ 内の cap の間の open annulus $U_0$ は $P\setminus K$ の
一つの連結成分であり，その全体が $D_1$ に含まれる．
\[
 \kappa_1=\int_{D_1\cap K}\omega,\qquad
 M_0=\int_{U_0}\omega>0
\]
と置く．$\operatorname{Diff}_0(P\operatorname{ rel}K)$ の中で foliation を変更しても，
$U_0$ 全体は第1 basin 内に残る．従って必ず
\[
 \operatorname{area}_{\omega}(D_1^{\rm new})\ge\kappa_1+M_0
\]
である．

ところが target $a_1=\kappa_1+M_0/2$ は個別の strict margin
$a_1>\kappa_1$ を満たす．元の第1 basin から減らした面積を第2 basin に足し，
他の target areas を元のままにすれば，どの target も正で，それぞれの
protected-area margin と全体の総量一致を満たす．しかし上の不変な下界に反する．
Parameter interval の両端では元の target，中央でこの target となるよう，
非負 bump function による凸補間を取れば，open end germs を保つ smooth data
に対しても同じ障害が生じる．
これは元の disk-complement 版の反例ではなく，追加 annulus を固定する版への
無条件な拡張の反例である．
\end{example}

\begin{remark}[Full critical neighborhoods と actual 適用の範囲]
すべての full critical leaves を含む pants neighborhoods $Q_v$ まで固定し，
各 $U_\alpha$ が一つの critical region にしか接しないなら，
column 条件はその region に元の面積を全て割り当てる．この場合 target critical
areas は元の値に強制され，面積を変更する自由はない．
既存の source-marking isotopy が $Q_v$ を固定する事実から，将来の
critical-area 調整でも $Q_v$ 全体を固定できるとは結論しない．

本定理は，protected part の外側で上記の smooth reduced data と relative marked
interpolation が実際に得られる場合の実構成である．Source 固定点を含む singular
reduction 全体に smooth area-form family を仮定していない．
固定した nodal neighborhoods の外側に適切な compact free reduction を取り，
その components，incidence，free-cell masses の smoothness を確認すればよい．
例えば $a_h(b)=\nu_b(D_h(b))+\delta a_h(b)$ とすると，
$r_h(b)=\sum_\alpha\int_{D_h(b)\cap U_\alpha}\omega_b+\delta a_h(b)$
なので，$\delta a_h$ の smoothness だけでよく，protected mass の非滑らかさは相殺される．
ただしこの事実は target の有限輸送条件を自動的に保証するものではない．
また foliation の面積を実現する定理であって，異なる既製 symplectic pieces の
primitive periods 又は core collar germs を自動的に一致させる定理ではない．
\end{remark}


\subsection{面積配分の有限 cut 判定}
\label{ng8:area-cut-section}

前定理の正の配分条件は，有限個の線型不等式で検査できる．
これは max-flow/min-cut の既知の原理
\cite[Theorem~1, p.~400]{FordFulkerson1956}に対応する有限二部問題であり，
一般の組合せ原理の新規性は主張しない．以下では必要な有向の配分問題を直接証明する．

\begin{proposition}[Strict な面積輸送の cut 条件]
\label{ng8:strict-area-cuts}
前定理の incidence graph を，region vertices $h$ と可動成分 vertices $\alpha$
を結ぶ二部 graph $G$ とする．辺は $(h,\alpha)\in E$ のとき，かつそのときに限り置く．
各非孤立 region の要求量を $r_h=a_h-\kappa_h>0$，各可動成分の面積を
$M_\alpha>0$ とする．孤立 region は $r_h=0$ を要求する．
可動成分 vertex は孤立しない．
全ての辺で正の配分が存在するための必要十分条件は，次の二つである．
\begin{enumerate}
\item 各 connected component $G_0$ の region 集合 $H_0$ と可動成分集合
  $A_0$ について，
  \[
    \sum_{h\in H_0}r_h=\sum_{\alpha\in A_0}M_\alpha .
  \]
\item 各 $G_0$ と各 $\varnothing\ne S\subsetneq H_0$ について，
  $N(S)\subset A_0$ を隣接する可動成分の集合とすると，
  \begin{equation}\label{ng8:eq:strict-area-cut}
    \sum_{h\in S}r_h<
    \sum_{\alpha\in N(S)}M_\alpha .
  \end{equation}
\end{enumerate}
従って検査する cut は有限個である．$r_h,M_\alpha$ が compact parameter
space 上で連続で，graph が固定され，上の条件が全 parameter で成り立つなら，
全ての strict inequalities と正値性には一様な正の余裕がある．
\end{proposition}

\begin{proof}
必要性を示す．Component 内の row sums と column sums を足せば最初の等式を得る．
正の配分 $x$ に対して
\[
 \sum_{\alpha\in N(S)}M_\alpha-\sum_{h\in S}r_h
 =\sum_{\substack{h\notin S,\ \alpha\in N(S)\\(h,\alpha)\in E}}x_{h\alpha}.
\]
Connected component の proper nonempty row subset には，右辺に少なくとも
一つの辺がある．従って差は正である．

十分性には，まず非負配分の判定を証明する．正の margins $r,M$ と
component-wise な総和一致を仮定し，上の全 cut inequalities を非strict に
弱める．非負 entry $y_{h\alpha}$ で row sums は $r_h$ 以下，column sums は
$M_\alpha$ 以下となるものの compact polytope 上で，総和を最大にする $y$ を取る．
最大値が要求総量より小さければ，未充足の row から始める．Row から隣接する
column へは増量方向に進め，column から row へは $y_{h\alpha}>0$ の辺だけを
減量方向に進める．未充足の column に到達できれば，有限の alternating path
に沿って十分小さい正量を増減し，総和を増やせて矛盾する．

従って到達した row 集合 $S$ の全隣接 columns $N(S)$ は充足されている．
しかも $S$ 外の row から $N(S)$ への正の entry はない．あれば逆向きの辺で
その row に到達するからである．最初の row は未充足なので，
\[
 \sum_{\alpha\in N(S)}M_\alpha
 =\sum_{h\in S,\alpha}y_{h\alpha}
 <\sum_{h\in S}r_h
\]
となり，非strict cut 条件に反する．従って非負の完全配分が存在する．

今度は全辺へ同じ十分小さい $\epsilon>0$ を先に配り，
\[
 r'_h=r_h-\epsilon\deg(h),\qquad
 M'_\alpha=M_\alpha-\epsilon\deg(\alpha)
\]
を残りの margins とする．Component-wise な総和一致は保たれる．
Cut の余裕の減少は
$\epsilon\,\#E(H_0\setminus S,N(S))$ だけなので，有限個の strict inequalities
と正値性から，残りの margins は非負配分の条件を満たすようにできる．
その非負解 $y$ に $\epsilon$ を加えれば，全辺で正の解になる．
最後の一様性は有限個の正の連続関数の compactness による．
\end{proof}

\begin{remark}[端だけの検査で足りる場合]
\label{ng8:affine-area-cuts}
一つの compact interval 上で graph が固定され，
全ての $r_h(\mu),M_\alpha(\mu)$ が affine であるなら，
component-wise な等式と strict cut inequalities は両端で検査すれば足りる．
各 cut の差も affine だからである．
Source を固定した actual 構成では，free 部の面積が smooth でも affine とは
限らず，この端点判定を無条件に用いてはならない．
\end{remark}

\begin{example}[両端の面積配分だけでは足りない]
\label{ng8:interior-area-cut-obstruction}
一つの external basin $D_1$ 内に保持 annulus を置き，その内側の可動成分
$U_0$ が $D_1$ だけに接し，外側の可動成分 $U_1$ が三つの basins
$D_1,D_2,D_3$ に接するとする．Incidence edges は
$(1,0),(1,1),(2,1),(3,1)$ である．
固定領域の面積 $\kappa_h$ を加えた target を $a_h=\kappa_h+r_h$ とし，
$\mu\in[-1,1]$ 上で
\[
 (r_1,r_2,r_3)=\left(\tfrac12,\tfrac34,\tfrac34\right),\qquad
 M_0(\mu)=\tfrac34-\tfrac12\mu^2,\qquad M_1(\mu)=2-M_0(\mu)
\]
と置く．全て正で，総和は常に一致する．両端では
$x_{10}=\tfrac14,x_{11}=\tfrac14,x_{21}=x_{31}=\tfrac34$ が正の配分である．
一方 $\mu=0$ では $x_{10}=M_0=\tfrac34>r_1=\tfrac12$ となり不可能である．
同じことは $S=\{2,3\}$ の cut
$r_2+r_3<M_1$ の不成立として読める．

この面積 family は，三 punctures の球面で $D_1$ の regular arm に固定 annulus
を置き，$U_0,U_1$ の内部で正の densities のみを変えて実現できる．
固定領域とその境界近傍の面積を先に小さく取り，その外に支持した bump を用いる．
例えば初期葉層の自由部分の cell masses を
\[
 m_{10}=\tfrac34-\tfrac12\mu^2,\quad m_{11}=\tfrac14,\quad
 m_{21}=m_{31}=\tfrac12+\tfrac14\mu^2
\]
と指定すれば，各 column の和は上の $M_\alpha$ となり，
初期葉層自身が両端で target を満たす．
さらに全ての $\mu^2$ を，両端の開近傍で $1$，中央の開近傍で $0$ となる
滑らかな $h(\mu)\in[0,1]$ に置き換えれば，open end germs を満たす模型にもなる．
これは元の actual KE 族の反例ではなく，
smooth free-band data と両端の確認だけから全 interval の面積実現を
推論できないことを示す模型である．
\end{example}


\subsection{Source の全近傍を保つ actual neck の微小面積指定}

以下では一つの許容された stage と $t\ne0$ を固定する．
Normal circle と canonical moment は既存 P0 のものを用いる．
面積と第二作用を比較する場合には周期 $2\pi$ に規格化するが，
moment の向きと加法定数を選び直さない．
本節の定数又は moment interval に $t\to0$ での一様性を要求しない．

\begin{lemma}[Source の飽和 annuli と滑らかな自由部分]
\label{ng8:p0-source-protected-reduction}
定理\ref{ke:projected-source-neck}で構成した actual neck と，
その後の end collars の相対整合を考える．
必要なら $0$ を含む compact moment interval $J$ を小さく取ると，
source を通る相異なる正則閉葉の小さい飽和 annuli
$A_1,\ldots,A_s$ を，critical leaves と保持する end collars から離して取れる．
Moment に沿う marked identification を用いた記法で，これらを固定 annuli とする．
このとき次が成り立つ．
\begin{enumerate}
\item $J\times A_a$ の actual circle orbit map による逆像は，
source の full compact pinched torus とその飽和近傍を含む．
これらの全 fiber partition を保持することは，source 点だけの固定より強い条件である．
\item 保持する end collars と全 $A_a$ を含む固定領域を $K$ とする．
$P\setminus K$ 上の actual normal circle は自由であり，
その reduced area forms は $\mu\in J$ に滑らかに依存する．
これは $\mu=0$ を跨いで成立する．
$K$ 内の reduced form，又は $K$ を含む領域の面積が
$\mu=0$ で滑らかであることは主張しない．
\end{enumerate}
\end{lemma}

\begin{proof}
補題\ref{ke:source-reduction-pinched-fiber}では orbit space は
$(\mu,p)\in I_t\times P$ で，非自由軌道は有限個の $(0,q_a)$ に限られる．
各 $q_a$ は相異なる閉 regular leaves 上にあり，saddle leaf と境界から離れている．
その小さい regular annular neighborhoods を互いに素に選ぶ．
定理\ref{ke:projected-source-neck}の基準葉層は残余変数 $p$ の葉層である．
End collars の整合は少し広い境界 collars 内に支持を持つので，
ここで選んだ内部 annuli を避けられる．
既に moment に依存する marked 葉族を用いている場合にも，
補題\ref{ke:simple-saddle-family-trivialization}で $J$ を縮め，
source markings を固定する reduced identification を選べる．
この identification を source 点上で ambient smooth map に持ち上げる必要はない．
最後の変更は source の annuli 上で恒等になるからである．

$\mu=0$ で source を通る閉葉の全逆像は pinched torus であり，
その両側の regular leaves と小さい moment interval を含む逆像は
その全 fiber の飽和近傍である．これが第一の結論である．
全 source は $K$ の内部に置かれる．従ってその外の閉じた自由部分では
moment map は正則で，$(\mu,p)$ は滑らかな submersion である．
Actual form の reduction と上の reduced identification により，
自由部分の area forms は $\mu=0$ を含めて滑らかである．
Source 上の orbit-space 座標又は reduced form の滑らかさを用いていない．
\end{proof}

\begin{corollary}[P0 の source を固定した微小 critical-area 指定]
\label{ng8:p0-small-area-flexibility}
前補題の固定 $t,J$ と marked identification を用いる．
保持領域 $K$ は end collars と source regular annuli の和とし，
補助的な puncture caps を用いる場合はそれらも固定する．
Critical leaves の全近傍を $K$ に追加することは，この系では要求しない．
Annuli と collars は，各 critical region $D_h$ に自由な開部分が残るように選ぶ．

元の actual area を $a_h^0(\mu)$ とする．
ある $\epsilon>0$ が存在し，滑らかな関数 $\delta a_h:J\to\RR$ が
\begin{equation}\label{ng8:eq:p0-small-area-increments}
 \sum_h\delta a_h(\mu)=0,\qquad
 \max_h\|\delta a_h\|_{C^0(J)}<\epsilon
\end{equation}
を満たし，$J$ の両端の指定した open intervals 上で全て零なら，
critical-region areas を
\[
 a_h(\mu)=a_h^0(\mu)+\delta a_h(\mu)
\]
とする actual Lagrangian neck fibration を構成できる．
同じ canonical moment，normal circle，marked tree，end collars と端の open germs，
全 source の上記 compact nodal neighborhoods を保つ．
正則・saddle・source fibers の型と全 fiber partition の接着条件も保つ．

$a_h^0$，$a_h$ 又は $\int_{D_h\cap K}\omega_{t,\mu}$ の滑らかさは不要である．
滑らかさを要求するのは増分 $\delta a_h$ と自由部分の area forms だけである．
\end{corollary}

\begin{proof}
$U_\alpha$ を自由部分 $P\setminus K$ の連結成分とする．
Critical leaves は面積零であり，各 $D_h$ と $U_\alpha$ の内部の交わりが
非空であるとき，二部グラフ $G$ に辺 $(h,\alpha)$ を置く．
このグラフは連結である．実際，boundary collars を除いた曲面は連結である．
二点を結ぶ path を，critical leaves と source annuli の境界に横断的に選ぶ．
Critical leaf を横切るときは同じ $U_\alpha$ に留まるので，
その両側の region vertices は $U_\alpha$ を介して結ばれる．
Source annulus を横切るときは一つの regular critical region $D_h$ 内に留まるので，
その両側の自由成分は $D_h$ を介して結ばれる．
これは $G$ 内の path を与える．Annuli を十分小さく選んだので，全 region vertex が現れる．

辺ごとの元の自由面積を
\[
 m_{h\alpha}(\mu)=
 \int_{D_h\cap U_\alpha}\omega_{t,\mu}>0
\]
と置く．Marked identification の下で積分領域は固定され，
その閉包は source を避ける自由部分に含まれるので，$m_{h\alpha}$ は滑らかである．
有限グラフと compact な $J$ より
\[
 m_*:=\min_{(h,\alpha)\in E(G),\,\mu\in J}m_{h\alpha}(\mu)>0.
\]

$G$ の spanning tree を一つ固定する．総和零の vector $v=(v_h)$ に対し，
tree の葉から順に消去することで，辺上の vector $L(v)$ を線形に選び，
\[
 \sum_\alpha L(v)_{h\alpha}=v_h,
 \qquad \sum_h L(v)_{h\alpha}=0
\]
とできる．これは region vertices の需要を $v_h$，自由成分 vertices の需要を零とする
tree flow であり，最後の頂点の式は $\sum_hv_h=0$ から従う．
$\|L(v)\|_\infty\le C_G\|v\|_\infty$ と取る．
各非空 cell の内部に，$K$ と critical leaves を避けた非負の smooth bump
$\varphi_{h\alpha,\mu}$ を選び，その cell における $\omega_{t,\mu}$ に関する積分を
$1$ と正規化する．Marked identification と $J$ の compactness により
$B:=\max\{1,\sup_{h,\alpha,\mu}\|\varphi_{h\alpha,\mu}\|_\infty\}<\infty$ であり，
異なる cell の bumps は互いに素な支持を持つ．
そこで
\[
 \epsilon=\frac{\min\{m_*,B^{-1}\}}{2\max\{1,C_G\}}
\]
と取れば
\[
 x_{h\alpha}(\mu)=m_{h\alpha}(\mu)+L(\delta a(\mu))_{h\alpha}>0
\]
である．各 column の和は元の $U_\alpha$ の面積に等しく，
各 row の和だけが $\delta a_h$ だけ変わる．
端の open intervals 上では $x=m$ である．

この小摂動の場合には density 自体を明示できる．自由部分上で
\[
 \rho_\mu=1+\sum_{(h,\alpha)\in E(G)}
 L(\delta a(\mu))_{h\alpha}\,\varphi_{h\alpha,\mu},
 \qquad \widehat\omega_\mu=\rho_\mu\omega_{t,\mu}
\]
と置くと，$\rho_\mu\ge1/2$，各 cell の積分は $x_{h\alpha}$ であり，
各 $U_\alpha$ の面積差は零である．Density は $K$ と critical leaves の近傍で $1$，
また $\delta a=0$ となる全 parameter 領域で恒等的に $1$ である．
補題\ref{ng8:componentwise-primitive}により差の primitive を
$P\setminus K$ に compact support を持つよう選び，その自由成分ごとに相対 Moser を行う．
これは定理\ref{ng8:componentwise-critical-area}の明示的な微小版である．
Source を含む $K$ 内の形式を変更又は滑らかに延長する必要はない．
得られる reduced diffeomorphisms は $K$ の開近傍および
$\delta a=0$ となる全 parameter 領域で恒等となる．
移された葉層では，$K$ 内の元の面積はそのままで，自由部分の面積だけが
$\delta a_h$ だけ変わる．これが所期の等式を与える．

最後に circle lift を取る．各 regular lifted leaf の接空間は normal generator と
$\ker d\mu$ 内の一方向で張られるから，$\iota_V\omega_t=-d\mu$ により Lagrangian である．
Saddle graph の lift は元と同じ isotropic strata を持ち，source の全近傍は恒等である．
相対 Moser isotopy は自由部分に compact support を持つので，
principal circle bundle の connection で水平に持ち上げ，$K$ 上で零に延長できる．
これにより moment と circle action を保つ ambient diffeomorphism が得られ，
葉の全 partition を移す．この diffeomorphism の symplectic 性は要求しない．
Lagrangian 性は上の reduction で直接確認済みである．
従って marked quotient，connected fibers と compact Hausdorff 性も保たれる．
全 end collars および moment の両端で全葉が変わらないので，
既存 corner・core との whole-fiber overlaps はそのまま接着する．
\end{proof}

\begin{remark}[局所自由度と残る atlas 条件]
\label{ng8:p0-area-scope}
この系は，固定した actual KE neck に対する非自明な相対変形を実際に与える．
有限個の neck に同様の変更を行っても，既存の corner と core を変更せず接着できる．
同じ normal moment と全保持領域を保つため，
命題\ref{ke:moment-relative-replacement}の基底値比較にも適合する．
ただし $\epsilon$ は固定 stage・$t$・保護領域に依存してよい．

全 critical actions を指定した affine functions へ変えるためには，
まず primitive と cycle framings を固定し，source の winding jump を含めて
それに対応する \emph{relative area increments} を求めなければならない．
その増分が上の小ささ又は一般の有限輸送条件を満たすか，
端の open germs と一致するかは別の問題である．
全 interval の absolute areas を affine と仮定することも結論することもない．
Source の固定領域に含まれる面積の非滑らかさは，元と新で同じ寄与として保持される．
従ってこの微小実現を，五次・六次の同じ射影の atlas と GH の同時実現，
任意の指定 critical actions の実現，又は退化 parameter に一様な面積予算とはしない．
\end{remark}

\subsection{球面 pants 木における一意な面積配分}
\label{ng8:spherical-tree-area}

本項では，系\ref{ng8:p0-small-area-flexibility}の保持集合を具体的に用いる．
$P$ は球面から有限個の puncture disks を除いた連結な compact surface であり，
全 critical leaves は一つずつの simple saddle を持つとする．
その connected-leaf quotient を $q:P\to T$ と書くと，$T$ は pants tree である．
Critical regions $D_h$ は $T$ の開辺 $e_h$ に対応する．
補助的な puncture caps を付ける場合も，その cap と end collar は固定する．

保持集合 $K$ は end collars と，critical leaves を避ける互いに素な閉飽和
regular annuli $A_1,\ldots,A_m$ の和とする．各 $A_s$ は一つの $D_h$ に含まれ，
$q(A_s)$ は $e_h$ 内部の閉 interval である．Annuli 同士，critical leaves，
および end collars の間には開 collar の余裕を残す．同じ辺上の保護 intervals が
重なる場合には，先にその和の各連結成分へまとめ，その後の個数を $m$ と数える．
Source の複数の保護円が球面上で入れ子になることは許す．

\begin{proposition}[Source-annulus incidence graph の木性]
\label{ng8:spherical-incidence-tree}
$P\setminus K$ の連結成分を $U_\alpha$，critical regions の個数を $H$ とし，
$D_h\cap U_\alpha\ne\varnothing$ のときだけ辺 $(h,\alpha)$ を置く二部グラフを
$G$ とする．このとき $G$ は木である．さらに，$D_h$ が $n_h$ 個の保護 annuli を
含むなら，$D_h\setminus K$ の $n_h+1$ 個の連結な free pieces は相異なる
$U_\alpha$ に属し，
\[
 \deg_G(h)=n_h+1,\qquad
 \#\{U_\alpha\}=m+1,\qquad
 |E(G)|=H+m=|V(G)|-1.
\]
したがって，一つの cell $D_h\cap U_\alpha$ 自体も連結である．
$P$ がこのような球面成分の非交和なら，成分ごとの同じ主張により forest を得る．
\end{proposition}

\begin{proof}
End collars を除く操作は tree の端を短くするだけで，その連結性を変えない．
一つの保護 annulus は，$T$ の辺内部の閉 interval を取り除く操作に対応する．
Tree には別の迂回路がないので，この操作は自由部分の連結成分をちょうど一つ増やす．
逆像の連結性は，connected fibers を持つ closed quotient $q$ に対する
補題\ref{topo:connected-preimages}から従う．ここでは open subtree の逆像にも
同補題を用いられる．従って保護 annuli を順に除くと自由成分数は $m+1$ となる．
球面上の各 regular circle が separator であることによっても同じことが分かる．

一つの開辺 $e_h$ に $n_h$ 個の保護 intervals があれば，残る部分は
$n_h+1$ 個の非空 intervals である．二つの異なる intervals が同じ自由成分内で
再会すると，それらを隔てる保護 interval を避けて $T$ 内で結ぶ path ができる．
これは tree の path の一意性に反する．各 interval の逆像は一つの regular
annulus 又はその端を切った部分であるから，必要な cell の連結性も従う．

$G$ の連結性も確認する．End collars の外の連結な曲面内で path を取り，
critical vertices を避け，critical leaves と保護 annuli の境界に横断的にする．
Critical leaf を横切るときは同じ $U_\alpha$ に留まるため，二つの region labels
はその $U_\alpha$ を介して結ばれる．保護 annulus を横切るときは一つの $D_h$ 内に
留まるため，両側の自由成分はその region vertex を介して結ばれる．
このように $G$ 内の path を得る．各 region と自由成分は非空 open set を含むので，
全 vertices がこの議論の対象となる．
最後に
\[
 |E(G)|=\sum_h(n_h+1)=H+m,\qquad |V(G)|=H+m+1
\]
である．連結グラフのこの等式から $G$ は木となる．
\end{proof}

\begin{corollary}[Edge cut による一意配分と全 target の判定]
\label{ng8:tree-edge-area-test}
前命題の保持集合を固定し，定理\ref{ng8:componentwise-critical-area}と同じ記法で
\[
 r_h=a_h-\kappa_h,\qquad M_\alpha=\nu(U_\alpha)>0
\]
と置く．$e=(h,\alpha)$ を取り除いた $G$ の二つの連結成分のうち，
region vertex $h$ を含む方を $C_e$ と書き，
\begin{equation}\label{ng8:eq:tree-edge-flow}
 X_e(a)=
 \sum_{k\in C_e\cap H}(a_k-\kappa_k)
 -\sum_{\beta\in C_e\cap U}M_\beta
\end{equation}
と定める．ここで $H,U$ は二種類の vertex 集合も表す．
所定 target を同じ marked relative orbit の中で実現できるための必要十分条件は
\begin{equation}\label{ng8:eq:tree-edge-inequalities}
 \sum_h(a_h-\kappa_h)=\sum_\alpha M_\alpha,
 \qquad X_e(a)>0\quad(e\in E(G))
\end{equation}
である．可解な場合の全 cell masses は一意であり，$x_{h\alpha}=X_{(h,\alpha)}(a)$
となる．これは mass allocation の一意性であって，density 又は foliation の一意性ではない．

Parameter family では，固定した incidence と smooth free data の下で，
$r_h,M_\alpha$ が smooth なら配分も上式により smooth である．
両端の target が元の area と一致する場合，配分もそこで元の cell masses と一致する．
従って同じ relative Moser theorem をそのまま適用できる．
$\kappa_h$ 個別の smoothness，source 内部での smooth reduced form，
又は新たな whole-fiber 仮定は必要ない．
\end{corollary}

\begin{proof}
各辺を region vertex から free-component vertex へ向け，vertex の signed balance を
$b(h)=r_h$，$b(\alpha)=-M_\alpha$ とする．
$C_e$ 内の vertex equations を足すと，内部辺は相殺し，境界に残る辺は $e$ 一つだけである．
その向きは $C_e$ の外向きなので，必ず $x_e=\sum_{v\in C_e}b(v)=X_e(a)$ となる．
これで一意性と必要性を得る．

逆に total balance を仮定して $x_e=X_e(a)$ と置く．
Region vertex $h$ を除くと，その各隣接辺の反対側の成分は全ての他の vertices を
分割する．Total balance により，それらの負の balances の和は $b(h)=r_h$ である．
従って $\sum_\alpha x_{h\alpha}=r_h$．
Free vertex $\alpha$ については，隣接する各 $C_e$ が他の vertices を分割するので，
$\sum_hx_{h\alpha}=-b(\alpha)=M_\alpha$ となる．
全 $X_e>0$ なら一般の正の有限配分条件を満たし，相対 Moser の構成が適用できる．
Parameter に関する主張は有限和の表示と同定理から従う．
\end{proof}

これは一般の strict Hall cut 判定のこの tree への特殊化である．
検査する不等式は $H+m$ 個だけでよく，一つの root を選んで subtree balances を
末端から加えることで全て求められる．Forest の場合は各 connected component に
一つずつ total balance を要求し，同じ edge formulas を用いる．
$r_h,M_\alpha$ が interval 上で affine なら端での検査で足りるが，
source 保持下でそれらが smooth であるだけの場合に端点検査へ置き換えることはできない．

\subsection{複数の保護 annuli と累積面積の制約}

元の cell masses を $m_e>0$，critical-region areas を $a_h^0$ とし，
$a_h=a_h^0+\delta a_h$，$\sum_h\delta a_h=0$ とする．前系の式は
\begin{equation}\label{ng8:eq:tree-increment-flow}
 x_e=m_e+\sum_{h\in C_e\cap H}\delta a_h
\end{equation}
となる．従って全可解増分は，零総和の hyperplane 内で，有限個の明示的な
strict linear inequalities によって切り出される開凸集合である．
この集合は零を含むが，任意の大きい増分を含むとは限らない．

同一 region $h$ に $n\ge1$ 個の annuli を辺に沿って並べる．
両端の free cells を $C_{h,-},C_{h,+}$，その元の面積を $m_{h,-},m_{h,+}$ とする．
中間の $n-1$ 個の cells は，それ自体が自由成分であり，$G$ では degree-one
free vertices になる．従ってその面積は，いかなる許容された変更でも元の値に固定される．
両端の cells は相異なる自由成分に属する．

元の pants tree から辺 $e_h$ の内部を除いて得る両側の region-edge 集合を
$H_h^-,H_h^+$ とする．これらは $h$ を含まず，
$H_h^-\coprod H_h^+=H\setminus\{h\}$ である．
$-$ 側の自由成分への incidence edge を切ったとき，$h$ 側に残る regions は
$\{h\}\cup H_h^+$ だから，\eqref{ng8:eq:tree-increment-flow}と零総和条件より
\begin{equation}\label{ng8:eq:two-end-cells}
 x_{h,-}=m_{h,-}-\sum_{k\in H_h^-}\delta a_k,
 \qquad
 x_{h,+}=m_{h,+}-\sum_{k\in H_h^+}\delta a_k.
\end{equation}
所定 target の制約は，これら二つの量の strict positivity である．
中間の cells については $x=M>0$ が自動的に成り立つ．
保護 annulus がない region では一つの free cell だけが存在し，
対応する条件は $a_h>\kappa_h$ となる．

\begin{corollary}[External basin の入れ子と internal annulus の二端条件]
\label{ng8:ordered-source-area-cuts}
External basin の保護 annuli を critical leaf 側から puncture 側へ並べる．
Critical leaf 側の最初の保護 annulus から puncture までの全固定 tail の面積を
$L_h$ とする．これは保護部分だけでなく，その向こう側に隔離された全自由成分の
面積を含む．この basin に対応する条件はちょうど
\[
 a_h>L_h
\]
である．より内側に annuli を追加しても，この同じ tail 内を細分するだけなら
region 総面積について新しい不等式は生じない．

Internal region では，二つの critical leaves に最も近い保護 annuli の外側にある
両端 cells が\eqref{ng8:eq:two-end-cells}の二つの条件を与える．
その間に annuli を追加しても，中間の面積がより細かく固定されるだけで，
region 総 target に対する条件は両端の条件のままである．
ただし異なる region 内の入れ子については，両側 subtree に含まれる
\emph{全ての} target 増分を式に入れなければならない．
個々の basin の面積下限だけで internal edge の条件を置き換えてはならない．
\end{corollary}

\begin{proof}
External edge の puncture 側には他の region edge がないので，
$H_h^+=\varnothing$ としてよい．従ってその側の cell と中間 cells の面積は全て固定され，
critical leaf 側の cell の面積だけが $\delta a_h$ だけ変わる．
その positivity は $a_h-L_h>0$ である．
Internal region では前述の二端公式を用いる．中間 cells は一つの $D_h$ だけに
属する自由成分なので，column mass 条件から不変である．
これで同じ region における複数 annuli と入れ子の主張を得る．
\end{proof}

\begin{remark}[幅を縮めても消えない面積予算]
Source を通る全 regular circle の近傍を固定する条件では，保護 annulus の幅を
零へ近づけても，その circle の向こう側に隔離された面積は一般に零にならない．
従って $\nu(K)$ を小さくすることだけでは大きい target の障害を除けない．
上の判定は，既存の actual reduced forms と同じ保持領域・marked identifications を
使うものであり，source の compact nodal neighborhoods 又は全 overlap の保持に
新しい条件を加えていない．中央の小さい moment interval に適用する場合と，
離れた anchor face に適用する場合は，それぞれの実際の $K$ と面積で別に判定する．
\end{remark}

\begin{example}[正 genus では incidence cycle が生じる]
\label{ng8:positive-genus-incidence-cycle}
二つの pair of pants を，それぞれの二つの境界を二本の annular tubes で結んで貼る．
残る二つの境界を end とする曲面は genus one，boundary components 二個である．
二つの simple saddles と radial ends を持つ葉層を取り，Reeb graph の二つの
saddle vertices の間には二本の平行な internal edges があるとする．
この二つの internal regions $D_1,D_2$ に一つずつ閉飽和 regular annulus を置き，
end collars とともに固定する．自由部分は二つの pants 側 $U_L,U_R$ に分かれる．
両 $D_1,D_2$ が両 $U_L,U_R$ に接するので，incidence graph は
\[
 h_1-U_L-h_2-U_R-h_1
\]
という四辺 cycle を含む．二つの external regions は対応する $U_L,U_R$ に
degree-one vertices として付く．この例の保護円はそれぞれ非分離的である．

例えば external regions の自由需要をそれぞれ $1/2$，internal regions の需要を
$r_1=r_2=1$，自由成分の質量を $M_L=M_R=3/2$ とする．
External entries を $1/2$ とすれば，internal entries は
\[
 \begin{pmatrix}x_{1L}&x_{1R}\\x_{2L}&x_{2R}\end{pmatrix}
 =\begin{pmatrix}s&1-s\\1-s&s\end{pmatrix},\qquad 0<s<1
\]
と自由に選べ，一意配分は成立しない．
これらの正の質量は各 open cell 内の density と小さい固定 collar masses を用いて
実現できる．従って球面又は pants tree の仮定を外して，edge-cut の一意公式を
一般の protected set へ拡張することはできない．
\end{example}

\subsection{Source を跨ぐ面積要求と作用の framing}

この項では $t\ne0$ と一つの stage を固定する．以下の規格化は
補題\ref{ke:source-reduction-pinched-fiber}および
補題\ref{ke:anchor-liouville-actions}と同じ\emph{周期 $1$}である．
Normal cycle を $\alpha$，その generator を $V$ とし，
一つの primitive $\lambda_N$ に対して
\[
 \mu=\lambda_N(V)=\int_\alpha\lambda_N,\qquad
 \iota_Vd\lambda_N=-d\mu
\]
とする．$\lambda_N$ は circle 不変に取る．
Actual canonical primitive からこの選択へは circle averaging 又は
補題\ref{ke:circle-exact-transport}の exact transport を用いられ，
全 closed-cycle periods は変わらない．Source では $V=0$ なので $\mu=0$ である．
前の周期 $2\pi$ による模型と比較する際の moment は $\mu/(2\pi)$ であって，
canonical moment の向き又は加法定数を選び直さない．

\begin{lemma}[二つの capping framing と source の面積]
\label{ng8:source-capping-periods}
標準模型 $\mathbb C^2$ において
\[
 I_j=|z_j|^2/2,\quad
 \lambda_0=I_1d\theta_1+I_2d\theta_2,\quad
 m=I_1-I_2,\quad \mu=2\pi m,
 \quad w=z_1z_2=re^{\sqrt{-1}\phi}
\]
と置く．角座標 $\theta_j,\phi$ は周期 $2\pi$ である．$w\ne0$ 上の
sections $\theta_1=0$ と $\theta_2=0$ が与える reduced primitives は
\begin{equation}\label{ng8:eq:source-two-primitives}
 b_+=\frac{\sqrt{m^2+r^2}-m}{2}\,d\phi,
 \qquad
 b_-=\frac{\sqrt{m^2+r^2}+m}{2}\,d\phi,
 \qquad b_--b_+=\frac{\mu}{2\pi}\,d\phi .
\end{equation}
$\mu>0$ では $b_+$，$\mu<0$ では $b_-$ が $w=0$ へ滑らかに延びる．
Source を避けて滑らかに動く正向き境界を持つ領域 $D_\mu$ について，
その境界の $0$ の周りの winding を $n$ とし，
$B_D(\mu)=\int_{\partial D_\mu}b_+$ と置くと
\begin{equation}\label{ng8:eq:source-area-wall}
 \operatorname{Area}(D_\mu)=B_D(\mu)+n\min\{\mu,0\},
 \qquad B_D\in C^\infty .
\end{equation}
通常の領域では $n$ は含まれる source の数であり，一つの source を
含む場合は $n=1$ である．したがって面積の右微分と左微分の差は $-n$ となる．

正向き $\phi$-loop 上の capping cycles を $\beta_+,\beta_-$ とすると
\begin{equation}\label{ng8:eq:source-capping-shear}
 \beta_-=\beta_++\alpha,
 \qquad
 \int_{\beta_-}\lambda_0=\int_{\beta_+}\lambda_0+\mu .
\end{equation}
Winding $n$ の境界では右辺の補正はそれぞれ $n\alpha,n\mu$ である．
\end{lemma}

\begin{proof}
$|w|^2=4I_1I_2$ より
$I_1=(\sqrt{m^2+r^2}+m)/2$，
$I_2=(\sqrt{m^2+r^2}-m)/2$ である．
各 section への $\lambda_0$ の引き戻しが
\eqref{ng8:eq:source-two-primitives}を与え，両者の微分は
\[
 \omega_{\mathrm{red},\mu}
 =\frac{r}{2\sqrt{m^2+r^2}}\,dr\wedge d\phi
\]
に等しい．正側では $I_2$，負側では $I_1$ が $r=0$ で零となり，
それぞれ対応する primitive が延びる．Stokes の定理により
正側の面積は $B_D$，負側の面積は $B_D+n\mu$ である．
零での密度は可積分であり，両側の連続極限が零 reduction の面積を与える．
境界は $r=0$ を避けるので $B_D$ は $\mu$ の滑らかな関数である．

$(\theta_1,\theta_2)$ の winding vector では
$\alpha=(1,-1)$，$\beta_+=(0,1)$，$\beta_-=(1,0)$ である．
したがって $\beta_-=\beta_++\alpha$ であり，period の式も従う．
ここで切り替えているのは source を\emph{埋める} section である．
Source から離れた同じ境界 torus family 上で $\beta_+$ を滑らかに運ぶことと，
正負両側で capping section を使うことを混同しない．
\end{proof}

\begin{lemma}[Actual source の面積の非滑らかな項]
\label{ng8:actual-source-area-wall}
補題\ref{ng8:p0-source-protected-reduction}の固定 actual neck を取り，
$0$ を含む moment interval $J$ を必要なら小さくする．
$D_{h,\mu}$ は marked identification で固定された critical region とし，
境界は source を避けるものとする．保護領域 $K$ は各 source の全近傍を含み，
$D_h$ の内部の source は全て $K\cap D_h$ の内部に含まれるとする．
Source を除いた境界と marked quotient は $\mu$ に滑らかに依存するものとする．
これは同補題の end collars，caps と source の飽和 annuli による選択で満たされる．

$D_h$ 内の source 数を $n_h$，元の面積と保護質量をそれぞれ
$a_h^0(\mu)$，$\kappa_h(\mu)$ と書く．零 reduction の source 点には
質量を与えず，正則部分上の reduced area を積分する．このとき
\begin{equation}\label{ng8:eq:actual-wall-smooth-remainders}
 a_h^0(\mu)=A_h(\mu)-\frac{n_h}{2}|\mu|,
 \qquad
 \kappa_h(\mu)=K_h(\mu)-\frac{n_h}{2}|\mu|,
 \qquad A_h,K_h\in C^\infty(J).
\end{equation}
特に $a_h^0-\kappa_h$ は滑らかである．定数と $J$ は固定 stage，$t$，
保護領域に依存してよい．
\end{lemma}

\begin{proof}
各 source における circle representation は weights $(1,-1)$ である．
まず，ここで必要な equivariant Darboux chart を局所的に作る．
不変 Riemannian metric の exponential chart で作用を線形化し，原点の
symplectic form を定数形式 $\omega_0$ とする．十分小さい不変 ball 上で
$\omega_s=\omega_0+s(\omega-\omega_0)$ は symplectic である．
Radial homotopy により $\omega-\omega_0$ の不変 primitive を取り，
これを用いる Moser vector field を原点で零に取る．その flow は equivariant であり，
縮小した近傍を標準 symplectic 表現へ同定する．
同定後の moment の微分は $2\pi d(I_1-I_2)$ で，原点の値も零なので
$\mu=2\pi(I_1-I_2)$ が厳密に成立する．
また $\lambda_N-\lambda_0$ はこの可縮 chart 上で exact である．
したがって局所 closed-cycle periods に追加の実数定数は生じない．
この構成は元の residual 座標 $p$ を保存するとは主張しない．

各 source の chart 内で，互いに素な不変 balls
$B_s=\{I_1+I_2\le c_s\}$ を $K\cap D_h$ の持ち上げの内部に取る．
$J$ を縮めて $|\mu|<2\pi c_s$ とすれば，その reduced area は
\[
 \operatorname{Area}((B_s)_{\mathrm{red},\mu})
 =\pi c_s-\frac12|\mu|
\]
である．これは補題\ref{ng8:source-dh-wall}の積分を現在の規格化へ直した式である．
Ball の境界はこの $J$ 上で moment の submersion となり，作用は自由である．

これらの ball の内部を除いた $D_h$ と $K\cap D_h$ の持ち上げは，
境界も含めて compact な free quotient family を与える．
元の marked quotient の境界と，Darboux ball の境界はここではともに
滑らかな free 部分にあるので，同時に trivialize できる．
その reduced forms と領域は滑らかに依存し，積分は滑らかな関数となる．
球の寄与を加えれば\eqref{ng8:eq:actual-wall-smooth-remainders}を得る．
Critical boundary が saddle graph である場合は，source から離れた
有限個の滑らかな arcs への分割，又は固定した相対 rounding で同じ積分を行う．
Source 点で marked quotient の通常の滑らかさを仮定してはいない．
\end{proof}

\begin{proposition}[Source を補正した target 面積要求]
\label{ng8:source-corrected-area-demand}
前補題の同じ $K$ と marked relative isotopy class を保つ．
目標の critical-region areas を $a_h^{\mathrm{tar}}$ とすると，
定理\ref{ng8:componentwise-critical-area}に入れる自由部分の需要
$r_h=a_h^{\mathrm{tar}}-\kappa_h$ が滑らかであるための必要十分条件は
\begin{equation}\label{ng8:eq:source-target-wall-condition}
 a_h^{\mathrm{tar}}(\mu)-n_h\min\{\mu,0\}\in C^\infty(J)
\end{equation}
である．同値に $a_h^{\mathrm{tar}}+n_h|\mu|/2$ が滑らかである．
例えば両側の目標が $c_h^\pm+p_h^\pm\mu$ である場合には
\[
 c_h^-=c_h^+,
 \qquad p_h^-=p_h^++n_h
\]
が必要十分である．$n_h>0$ の領域の\emph{面積自体}を一つの affine function に
指定することは，この滑らかな相対問題の条件を満たさない．

一方，critical actions を指定するときには，保持 boundary cycles を含む
同じ primitive と同じ transported cycle framings で Stokes の式を書く．
$\partial D_h$ の向き付き cycle 表示に現れる作用を $J_\ell$，その符号付き
係数を $e_{h\ell}\in\mathbb Z$ とすれば，ある固定整数 $l_h$ により
\begin{equation}\label{ng8:eq:source-boundary-stokes}
 a_h(\mu)=\sum_\ell e_{h\ell}J_\ell(\mu)+l_h\mu
                       +n_h\min\{\mu,0\}
\end{equation}
と書ける．Cycle の選択を変えれば $l_h$ も対応して変わるので，
その整数を省略した unmarked な式として使わない．
同じ marked relative class の元と目標を比較すると
\begin{equation}\label{ng8:eq:source-relative-actions}
 a_h^{\mathrm{tar}}-a_h^0
 =\sum_\ell e_{h\ell}(J_\ell^{\mathrm{tar}}-J_\ell^0).
\end{equation}
したがって source から離れた critical cycles において，
共通の transported frame で滑らかな作用要求を与えれば，相対面積要求も滑らかである．
特にその frame で affine な critical actions を求めることは，
面積自体を affine とすることとは異なる．
\end{proposition}

\begin{proof}
最初の必要十分性は\eqref{ng8:eq:actual-wall-smooth-remainders}から直ちに従う．
Piecewise affine の場合は，$n_h\min\{\mu,0\}$ を引いた二本の直線の
値と傾きを零で合わせればよい．

Stokes の式を確認する．各 source の小 ball を除いた領域を $W_h$ とする．
$W_h$ は境界付き曲面であり，$J$ は区間なので，その上の principal circle
bundle は自明である．一つの section を固定し，周期 $1$ の circle 座標を
$a\in\mathbb R/\mathbb Z$ と書けば $\lambda_N-\mu\,da$ は basic である．
各 level でその微分は reduced form なので，$W_h$ 上で Stokes を適用できる．
各内境界の section cycle と局所の $\beta_+$ の差は整数倍の $\alpha$ であり，
その period 差は整数倍の $\mu$ である．
外境界の section cycles と，指定された critical／保持 cycles の差も同様である．
各小 ball の面積には\eqref{ng8:eq:source-area-wall}を使う．
内境界の smooth period は相殺し，残る整数倍を $l_h\mu$ とまとめると
\eqref{ng8:eq:source-boundary-stokes}を得る．
Saddle graph の境界 cycle についても，free 部分の有限 arcs 又は近傍の
regular cycles に Stokes を適用して極限を取れば同じ式になる．

相対 isotopy は source の保持近傍と boundary markings を固定する．
したがって各整数 winding と source 数は変わらず，同じ $l_h,n_h,e_{h\ell}$ を使える．
二つの式の差が\eqref{ng8:eq:source-relative-actions}である．
ここでは $\lambda_N$ が全四次元領域で closed であるとはしていない．
Homology による cycle の比較は同じ Lagrangian torus 又はその isotropic strata 上で行い，
領域間の比較には上の Stokes の式を使っている．
\end{proof}

\begin{remark}[この条件と有限面積輸送の関係]
\eqref{ng8:eq:source-relative-actions}で得られる需要が滑らかであることだけでは，
その実現は結論しない．自由成分ごとの総量，正の有限配分と strict cuts，
保持する open end germs をさらに照合する．これらが定理\ref{ng8:componentwise-critical-area}
の仮定を満たせば相対面積構成を適用でき，小さい増分の場合は
系\ref{ng8:p0-small-area-flexibility}の明示的な構成を使える．
Critical action の全線形関係と，固定された boundary actions を反映した
\eqref{ng8:eq:source-boundary-stokes}を先に作る必要がある．
Source local normal form だけから，任意の affine targets 又は実際の P0 の
大きい所定 targets がこれらを満たすとは結論しない．
\end{remark}

\begin{proposition}[Source shear を含む二つの anchor の比較条件]
\label{ng8:source-corrected-anchor-matching}
二つの core--neck collars を $s=L,R$ とし，同じ neck primitive $\lambda_N$ と
同じ canonical normal action $\mu$ へ比較する．各 collar の記号を
\[
 q_s=\sigma_s\mu,\quad \sigma_s\in\{1,-1\},\qquad
 M_s=\delta_s-\rho_s<0
\]
とし，core primitive の anchor residual period を $A_s$ とする．
Core と neck の primitive 差の residual period を
$\kappa_s^{\mathrm{prim}}$ と書く．これは保護\emph{面積} $\kappa_h$ とは別の量である．
Core cycle に対する比較 framing を
$\beta_{N,s}=\varepsilon_s\beta_{\mathrm{core},s}+k_s\alpha_s$ とすると，
補題\ref{ke:anchor-liouville-actions}の後の比較式は
\begin{equation}\label{ng8:eq:source-anchor-local}
 p_s=k_s+\varepsilon_s\frac{A_s}{M_s},\qquad
 c_s=\varepsilon_s\left(\kappa_s^{\mathrm{prim}}
                              -\rho_s\frac{A_s}{M_s}\right).
\end{equation}
ここで $c_s+p_sq_s$ が core 側の指定 affine germ である．

指定された実際の regular torus の同定と cycle transport に関して
\[
 \beta_{N,L}=\varepsilon\beta_{N,R}+n\alpha,
 \qquad \varepsilon\in\{1,-1\},\quad n\in\mathbb Z
\]
が成り立つとする．この同定で二つの affine targets が一致するための
必要十分条件は
\begin{equation}\label{ng8:eq:source-anchor-matching}
 c_L=\varepsilon c_R,
 \qquad \sigma_Lp_L=\varepsilon\sigma_Rp_R+n .
\end{equation}
従って anchor data による条件は
\begin{align*}
 \varepsilon_L\left(\kappa_L^{\mathrm{prim}}-\rho_L A_L/M_L\right)
 &=\varepsilon\varepsilon_R
   \left(\kappa_R^{\mathrm{prim}}-\rho_R A_R/M_R\right),\\
 \sigma_L\left(k_L+\varepsilon_L A_L/M_L\right)
 &=\varepsilon\sigma_R\left(k_R+\varepsilon_R A_R/M_R\right)+n .
\end{align*}
これは period の比較条件であり，全 open overlaps 上の葉の一致又は
全域 atlas の実構成までを結論しない．
\end{proposition}

\begin{proof}
\eqref{ng8:eq:source-anchor-local}は既存の anchor 計算における
$c_s+\rho_s(p_s-k_s)=\varepsilon_s\kappa_s^{\mathrm{prim}}$
を書き直したものである．同じ torus 上では $\lambda_N$ の制限は closed なので
\[
 J_L=\varepsilon J_R+n\mu .
\]
$q_s=\sigma_s\mu$ を代入して affine functions の定数項と一次係数を比較すると，
\eqref{ng8:eq:source-anchor-matching}とその同値な二式を得る．
Local source の capping framing の交換は
\eqref{ng8:eq:source-capping-shear}によりこの整数 shear に寄与する．
Source の canonical moment は零なので，この交換だけから実数の切片補正は生じない．
\end{proof}

\begin{remark}[測るべき整数と実数 period]
上の $n$ は指定された実際の cycle transport による整数である．
一つの境界が source を $n$ 回囲み，正負の capping framings を比較する
局所模型では source の winding がその整数になる．一般の neck の両端を比較するときは，
triple 側の既定 framing と transport path も含めて計算しなければならない．
単に moment interval が零を跨ぐことから，任意の transported critical cycle の
framing が source 数だけ変わるとはしない．

二つの neck primitives が異なり，その closed 差の normal period が零なら，
その residual period $C$ を明記し，切片条件を
$c_L=\varepsilon c_R+C$ とする．一つの $\lambda_N$ を使う上の設定では $C=0$ である．
Node germ を固定するだけではこの residual constant は決まらない
（例\ref{ke:focus-local-primitive-shift}）．
Source shear は整数勾配の比較を変えるが，実数切片の不一致を修復しない．
実際の $\kappa_s^{\mathrm{prim}}$ と global framing が全て測定済みであるとは主張しない．
\end{remark}
